The \code{BPatch\_thread} class represents and controls a thread of
execution that is running in a process.

\begin{apient}
  void getCallStack(std::vector<BPatch_frame>& stack)
\end{apient}

\apidesc{
  This function fills the given vector with current information about
  the call stack of the thread. Each stack frame is
  represented by a  \code{BPatch\_frame}.
}

\begin{apient}
  dynthread_t getTid()
\end{apient}

\apidesc{
  This function returns a platform-specific identifier for this
  thread. This is the identifier that is used by the
  threading library. For example, on pthread applications this
  function will return the thread\'s \code{pthread\_t} value.
}

\begin{apient}
  Dyninst::LWP getLWP()
\end{apient}

\apidesc{
  This function returns a platform-specific identifier that the
  operating system uses to identify this thread. For
  example, on UNIX platforms this returns the LWP id. On Windows this
  returns a HANDLE object for the thread.
}

\begin{apient}
  unsigned getBPatchID()
\end{apient}

\apidesc{
  This function returns a Dyninst-specific identifier for this
  thread. These ID's apply only to
  running threads, the BPatch ID of an already terminated thread my
  be repeated in a new thread.
}

\begin{apient}
  BPatch_function *getInitialFunc()
\end{apient}

\apidesc{
  Return the function that was used by the application to start this
  thread. For example, on pthread applications this
  will return the initial function that was passed to \code{pthread\_create}.
}

\begin{apient}
  unsigned long getStackTopAddr()
\end{apient}

\apidesc{
  Returns the base address for this thread's stack.
}

\begin{apient}
  bool isDeadOnArrival()
\end{apient}

\apidesc{
  This function returns true if this thread terminated execution
  before Dyninst was able to
  attach to it. Since Dyninst performs new thread detection
  asynchronously, it is possible
  for a thread to be created and destroyed before Dyninst can attach
  to it. When this hap-
  pens, a new \code{BPatch\_thread} is created, but isDeadOnArrival
  always returns true for this
  thread. It is illegal to perform any thread-level operations on a
  dead on arrival thread.
}

\begin{apient}
  BPatch_process *getProcess()
\end{apient}

\apidesc{
  Return the  \code{BPatch\_process}  that contains this thread.
}

\begin{apient}
  void *oneTimeCode(const BPatch_snippet &expr, bool *err = NULL)
\end{apient}

\apidesc{
  Cause the snippet expr to be evaluated by the process immediately.
  This is similar to the  \code{BPatch\_process::oneTimeCode}
  function, except that the snippet is guaranteed to run only on this
  thread. The process must be stopped to call this
  function. The behavior is synchronous; oneTimeCode will not return
  until after the snippet has been run in the
  application.
}

\begin{apient}
  bool oneTimeCodeAsync(const BPatch_snippet  &expr, void *userData = NULL,
                        BpatchOneTimeCodeCallback cb = NULL)
\end{apient}

\apidesc{
  This function sets up the snippet expr to be evaluated by this
  thread at the next available opportunity. When the
  snippet finishes running, Dyninst will callback any function
  registered through BPatch::registerOneTimeCodeCallback,
  with userData passed as a parameter. This function returns true if
  expr was posted or false otherwise.
  This is similar to the  \code{BPatch\_process::oneTimeCodeAsync}
  function, except that the snippet is guaranteed to run only on
  this thread. The process must be stopped to call this function. The
  behavior is asynchronous; oneTimeCodeAsync
  returns before the snippet is executed.
}
